\documentclass[12pt]{article}

\usepackage{sbc-template}
\usepackage{graphicx,url}
\usepackage[utf8]{inputenc}
\usepackage[brazil]{babel}
\usepackage{amsmath}
\usepackage{booktabs}
%\usepackage[latin1]{inputenc}  

     
\sloppy

\title{RS}

\author{Luciana P. Nedel\inst{1}, Rafael H. Bordini\inst{2}, Flávio Rech
  Wagner\inst{1}, Jomi F. Hübner\inst{3} }


\address{Instituto de Informática -- Universidade Federal do Rio Grande do Sul
  (UFRGS)\\
  Caixa Postal 15.064 -- 91.501-970 -- Porto Alegre -- RS -- Brazil
\nextinstitute
  Department of Computer Science -- University of Durham\\
  Durham, U.K.
\nextinstitute
  Departamento de Sistemas e Computação\\
  Universidade Regional de Blumenal (FURB) -- Blumenau, SC -- Brazil
  \email{\{nedel,flavio\}@inf.ufrgs.br, R.Bordini@durham.ac.uk,
  jomi@inf.furb.br}
}

\begin{document} 

\maketitle

\begin{abstract}
  This meta-paper describes.
\end{abstract}
     
\begin{resumo} 
  Este meta-artigo descreve.
\end{resumo}


\section{Introducción} \label{sec:intro}

%Sugestão: 1ºparágrafo: comecem dizendo qual é a motivação para se investigar RS.
La Inteligencia Artificial Neuro-Simbólica (NeSy) es un enfoque que busca combinar la capacidad de aprendizaje de las redes neuronales con el razonamiento estructurado de la IA simbólica. De esta manera, aspira a desarrollar sistemas más robustos, interpretables y capaces de generalizar mediante la incorporación de conocimiento previo \cite{Colelough2024, Yang2025, Acharya2025}. Tradicionalmente, se ha asumido que la integración de restricciones simbólicas guía el aprendizaje hacia representaciones conceptuales correctas. Sin embargo, estudios recientes \cite{Marconato2023-NotAll, Yang2024} han demostrado que los modelos NeSy pueden alcanzar un alto rendimiento predictivo apoyándose en conceptos cuya semántica no se corresponde con la esperada. Siguiendo la terminología propuesta por \cite{Marconato2023-NotAll}, este fenómeno se conoce como \textit{Reasoning Shortcuts} (RSs) y constituye un problema crítico que compromete la interpretabilidad, la fiabilidad y la capacidad de generalización fuera de distribución de estos sistemas. En este contexto, comprender y cuantificar la presencia de RSs resulta fundamental para evaluar de manera crítica las verdaderas capacidades de los modelos NeSy.


%2º parágrafo: definições anteriores de RS e quais as limitações delas.
El fenómemo de los \textit{reasoning shortcuts} no se limita a modelos NeSy. En el contexto de Deep Learning, \cite{Geirhos2020} describen atajos de aprendizaje como reglas de decisión que se desempeñan bien en benchmarks estándar, pero que fallan en condiciones de prueba más exigentes, como escenarios del mundo real.
En el ámbito neuro-simbólico, \cite{Li2024} introducen el término \textit{shortcut satisfaction}, donde el modelo se sobreajusta a una asignación de conceptos particular que satisface las restricciones lógicas impuestas. Sin embargo, múltiples configuraciones de conceptos pueden satisfacer las restricciones pero no todas corresponden a la verdad. Por otra parte, \cite{Marconato2023-NotAll} proponen una formalización probabilística de los RS, modelándolos como óptimos no intencionados en un proceso generativo consistente con el conocimiento previo $K$. Bajo este marco, un modelo incurre en RS cuando aprende una distribución sobre conceptos $p_\theta(C \mid X)$ que maximiza la verosimilitud en los datos de entrenamiento pero no coincide con la distribución de conceptos \textit{ground-truth} $p^*(G \mid X)$. A pesar de estos avances, las definiciones existentes presentan limitaciones importantes: se centran en propiedades del proceso de aprendizaje sin ofrecer mecanismos prácticos para cuantificar la presencia de atajos en modelos entrenados. Incluso las formulaciones más recientes carecen de métricas operativas que permitan medir de forma directa hasta qué punto las predicciones correctas de un modelo están respaldadas por conceptos semánticamente válidos.

%3º parágrafo: descrever a proposta de metricas do artigo e como elas resolvem as limitações anteriores.
Para abordar estas limitaciones, este trabajo propone una métrica diseñada para cuantificar el uso de \textit{reasoning shortcuts} a nivel de modelo. A diferencia de enfoques previos, que analizan los RS desde una perspectiva global de la tarea \cite{Marconato2023-NotAll} o mediante métricas indirectas sobre el espacio de conceptos como \textit{concept collapse} \cite{Bortolotti2024}, esta propuesta se centra en evaluar directamente en qué medida las predicciones correctas del modelo están respaldadas por conceptos semánticamente correctos. En particular, introducimos el \textit{Reasoning Shortcut Ratio} (RSR), que mide la proporción de aciertos que pueden atribuirse a configuraciones conceptuales incorrectas, proporcionando así una estimación interpretable del grado en que el modelo depende de shortcuts. Adicionalmente, extendemos esta métrica para considerar la ambigüedad estructural inherente al problema, ponderando las instancias según el número de configuraciones conceptuales que pueden explicar una misma etiqueta. Finalmente, proponemos una extensión probabilística que permite capturar no solo la presencia de shortcuts, sino también su intensidad, aprovechando la distribución de probabilidad sobre conceptos producida por modelos NeSy modernos. En conjunto, estas métricas permiten: $i)$ detectar el uso de reasoning shortcuts en modelos NeSy, $ii)$ comparar modelos más allá de métricas tradicionales como \textit{accuracy}, y $iii)$ analizar el impacto de la ambigüedad conceptual en el aprendizaje.

%4º parágrafo: dizer que aplicam as métricas em AES do ENEM, por ser um problema do mundo real, ao contrário dos toys examples comuns da literatura. Também dizer como aplicam as métricas e quais os experimentos feitos.
Para validar empíricamente las métricas propuestas, las aplicamos en el contexto de \textit{Automatic Essay Scoring} (AES) para el examen nacional de ingreso a la universidad --- ENEM ---, un problema del mundo real con alta relevancia social y complejidad semántica, en contraste con los \textit{toy examples} comúnmente utilizados en la literatura NeSy. Evaluamos nuestras métricas bajo diferentes configuraciones de aprendizaje, incluyendo entrenamiento supervisado y \textit{distant learning}, así como en distintos datasets: JBCS \cite{propor2024} que se componen de ensayos enviados por estudiantes a sitios web que simulan el examen ENEM y son evaluados por correctores del sitio web o por profesionales contratados; y Subcriteria \cite{feedbackpropor}, que proporciona anotaciones explícitas de conceptos en forma de subcriterios derivados de las rúbricas oficiales. Estos experimentos nos permiten analizar cómo varía el uso de reasoning shortcuts según el régimen de entrenamiento y la disponibilidad de supervisión conceptual. Los experimentos realizados permiten no solo cuantificar la presencia de RS en un escenario realista, sino también estudiar su relación con el desempeño predictivo y la calidad de las representaciones conceptuales aprendidas.

%5º parágrafo: dizer quais foram as conclusões e contribuições do trabalho.
% in progress...

\section{Trabajos Relacionados} \label{sec:relatedwork}

\paragraph{Neuroprobabilistic Logic Programming}  Many tools have been developed recently that enable neuroprobabilistic logic programming through weighted model counting. Examples of such are DeeProbLog \cite{DeepProbLog-Neurips,DeepProbLog-AI}, NeurASP \cite{NeurASP}, DPASP \cite{DPASP-Arxiv,DPASP-KR}, and Scallop \cite{Scallop-NEURIPS,Scallop-ACM}. Through these tools, it is possible to easily combine low level observations, such as classifications done by neural networks when receiving an input, with higher-order thinking done by logic formulae over the possible classifications. 

One other advantage is enabling two different types of training. In the first type, using only traditional supervised learning, it is possible to train the neural networks in what they are supposed to do, and then combine them through the logical formulae. The second type is called distant learning, where only the result of the inference is presented, and the gradient flows to the logical combinations that satisfy the inference. In this work, we employ both learning methods in order to assess our metrics. 

\paragraph{Reasoning Shortcuts Metrics}
Diversos trabajos han propuesto métricas para caracterizar la presencia de \textit{reasoning shortcuts} desde distintas perspectivas. Por un lado, \cite{Bortolotti2024, Marconato2023-NotAll} introducen el análisis de \textit{óptimos deterministas}, que cuantifica el número de soluciones conceptuales que satisfacen las restricciones lógicas de una tarea. Esta medida proporciona una indicación de la ambigüedad estructural del problema y, por ende, de su susceptibilidad a RS, aunque no permite evaluar directamente el comportamiento de un modelo entrenado. En el mismo trabajo, se propone la métrica de \textit{concept collapse}, que busca medir en qué medida las representaciones aprendidas combinan múltiples conceptos ground-truth. Un alto colapso sugiere que el modelo utiliza un conjunto reducido de conceptos para resolver la tarea; sin embargo, un bajo colapso no implica necesariamente la ausencia de RS, ya que el modelo podría estar activando múltiples conceptos de manera incorrecta sin reflejar la semántica adecuada. 

Por otro lado, en escenarios donde no se dispone de anotaciones de conceptos, \cite{Marconato2024-bears} propone un enfoque basado en la comparación entre múltiples modelos NeSy que alcanzan un alto desempeño en la tarea. La idea es analizar la consistencia entre los conceptos inferidos por estos modelos: discrepancias significativas pueden indicar la presencia de shortcuts. No obstante, este enfoque es indirecto, depende del entrenamiento de múltiples modelos y no proporciona una medida cuantitativa clara de la intensidad o frecuencia de los RS en un modelo individual.

\paragraph{RS Mitigation Strategies} Multiples estrategias han sido propuestas para mitigar la aparición de \textit{reasoning shortcuts} en modelos NeSy, las cuales pueden agruparse en enfoques supervisados, no supervisados y basados en el conocimiento. Entre las estrategias supervisadas, la \textit{concept supervision} \cite{Marconato2023-NotAll} introduce términos adicionales de pérdida para alinear las predicciones de conceptos con anotaciones \textit{ground-truth}. Si bien es efectiva, su aplicabilidad está limitada por el alto costo de obtención de anotaciones, así como por el riesgo de sobreajuste cuando estas cubren exhaustivamente el espacio conceptual. En una línea similar, el \textit{multi-task learning} \cite{Marconato2023-NotAll} busca reducir la ambigüedad entrenando modelos sobre múltiples tareas que comparten conceptos, aunque requiere disponibilidad de múltiples etiquetas y un diseño cuidadoso de tareas compatibles.

Por otro lado, los enfoques basados en el conocimiento, como \textit{knowledge editing} \cite{Marconato2023-NotAll}, intentan reducir la ambigüedad modificando las restricciones lógicas, pero su viabilidad depende de la capacidad de intervenir sobre el conocimiento previo. Asimismo, métodos como la \textit{abductive weak supervision} \cite{Yang2024} generan pseudo-etiquetas de conceptos mediante abducción lógica, aunque pueden introducir ruido en el proceso de aprendizaje.

Finalmente, diversas estrategias no supervisadas buscan regular el espacio de representaciones aprendidas. Entre ellas, la maximización de entropía \cite{Marconato2023-NotAll} y el \textit{concept smoothing} \cite{Yang2024} promueven distribuciones de conceptos menos deterministas, reduciendo el riesgo de ciertos tipos de RS, aunque sin garantías de alineación semántica. Métodos como la reconstrucción fomentan representaciones más informativas, pero introducen complejidad computacional adicional y dependen de supuestos sobre la estructura de los datos. Por su parte, el \textit{architectural disentanglement} \cite{Marconato2023-NotAll} restringe la expresividad del modelo procesando componentes independientes de la entrada, lo que puede reducir el espacio de shortcuts posibles, aunque requiere supuestos fuertes sobre la estructura de los datos. Adicionalmente, \cite{Stein2025} estudia los pitfalls de la programación neuro-simbólica, en donde 

En resumen, no existe una estrategia universalmente óptima. Los métodos supervisados tienden a ser más efectivos pero costosos, mientras que los no supervisados son más escalables pero menos fiables. En la práctica, la elección depende de las características de la tarea, la disponibilidad de anotaciones y el tipo de shortcuts presentes, siendo común la combinación de múltiples estrategias, aunque no siempre compatibles entre sí. En este trabajo, aprovec

\paragraph{Automatic Essay Scoring} (AES) is a field that aims to develop systems that are able to assess written text with minimal human intervention \cite{Page-1966}. In Brazil, given the social importance of the national university entrance exam --- the ENEM --- different datasets have been proposed in order to create such systems \cite{amorim-veloso-2017,Essay-br,ExtendedEssay-br,propor2024,paper-lais}. These datasets, except for the last, are composed of essays sent by students to websites that simulate the ENEM exam and are evaluated either by graders of the website or by hired professionals. The last one is composed by official ENEM essays with their respective official grades.

As the ENEM is a standardized exam with well defined, publicly available rubrics, the Subcriteria dataset \cite{feedbackpropor} was created by having essays annotated with their boolean properties as defined in the rubrics. These properties were called subcriteria, as they are used to define the criteria that compose each grade. After creating the dataset, the authors used Scallop in order to create a neuro-symbolic AES system that combines LLMs --- that ground the essays into the subcriteria --- with the logical formulae that define each grade. 
Although it has been shown that this neuro-symbolic system is competitive with the benchmarked SOTA \cite{JBCS} and is capable of explaining the reason for its grading, the authors only evaluated the explanations using agreement metrics. In the present work, we use this AES system and dataset in order to test and discuss our metrics that are specific for neuroprobabilistic logic programs.  


\section{Metodología} \label{sec:methods}

\subsection{Formalización de Reasoning Shortcuts}

Siguiendo la formulación de \cite{Bortolotti2024} y \cite{Marconato2023-NotAll}, consideramos tareas que requieren aprendizaje y razonamiento, las cuales predicen una etiqueta $\mathbf{y}$ a partir de una entrada de bajo nivel $\mathbf{x}$. Se asume que existe un conjunto de $k$ conceptos $\mathbf{c}^*$ de alto nivel asiciados a $\mathbf{x}$; entonces, $\mathbf{c}^*$ junto con el conocimiento previo $\mathrm{K}$, permiten inferir la etiqueta correcta $\mathbf{y}^*$ a través de una función $\beta_K: \mathbf{c^*} \mapsto \mathbf{y}$.

Un modelo neuro-simbólico puede representarse como una composición $\beta_K \circ \alpha$, donde $\alpha: \mathbf{c}^* \mapsto \mathbf{c}$ mapea los conceptos verdaderos a conceptos válidos para la predición de etiquetas. Idealmente, $\alpha$ debería recuperar los conceptos semánticamente correctos $\mathbf{c}^*$. Sin embargo, pueden existir múltiples funciones $\alpha$ que, incluso produciendo conceptos incorrectos, generan la etiqueta correcta al ser procesadas por $\beta_K$. Estas soluciones alternativas constituyen \textit{reasoning shortcuts} (RS).

Formalmente, un RS ocurre cuando existe una función $\alpha$ tal que:
\begin{equation*}
(\beta_K \circ \alpha)(\mathbf{c}) = \mathbf{y}^* \quad \wedge \quad \alpha(\mathbf{c}) \neq \mathbf{c}^*
\end{equation*}

Siguiendo \cite{Marconato2023-NotAll}, el número de óptimos deterministas (\textit{det-opts}) corresponde al número de funciones $\alpha$ que satisfacen esta condición para todos los ejemplos del conjunto de entrenamiento. Este valor puede expresarse como:

\begin{equation}
\text{det\_opts} = \sum_{\alpha\in\mathcal{A}} \mathbf{1}\left\{\bigwedge_{\mathbf{c}\in \text{supp}(\mathbf{c^*})}(\beta_K \circ \alpha)(\mathbf{c}) = \beta_k(\mathbf{c})\right\} = \prod_{\mathbf{y} \in \mathcal{Y}} |\mathcal{C}_{\mathbf{y}}|^{|\mathcal{C}_{\mathbf{y}}|}
\end{equation}

donde $\mathcal{C}_{\mathbf{y}}$ es el conjunto de configuraciones de conceptos que producen la etiqueta $\mathbf{y}$. Esta formulación muestra que, incluso en escenarios ideales, pueden existir múltiples soluciones óptimas que no corresponden a la semántica correcta, evidenciando la ambigüedad inherente en tareas NeSy.

En escenarios reales, no es posible enumerar explícitamente todas las funciones $\alpha$. En su lugar, se analiza el comportamiento del modelo entrenado $\alpha_{\text{modelo}}$ sobre un conjunto de datos $\mathcal{D} = \{(\mathbf{x}_i, \mathbf{c}_i, \mathbf{y}_i)\}_{i=1}^N$, identificando evidencia de RS cuando:
\begin{equation}
\hat{\mathbf{y}}_i = \mathbf{y}_i \quad \wedge \quad \hat{\mathbf{c}}_i \neq \mathbf{c}_i
\end{equation}

\subsection{Definición de la Métrica Propuesta}

A partir de la formalización anterior y motivados por las limitaciones de métricas existentes como el \textit{concept collapse} \cite{Bortolotti2024}, proponemos una métrica que cuantifica la proporción de predicciones correctas que son potencialmente explicadas por reasoning shortcuts.

Sea $\mathcal{D}$ un conjunto de $N$ instancias. Definimos el conjunto de aciertos del modelo como:
\begin{equation}
\mathcal{A} = \{i \mid \hat{y}_i = y_i\}
\end{equation}

y el subconjunto de aciertos que presentan evidencia de RS como:
\begin{equation}
\mathcal{A}_{RS} = \{i \mid \hat{y}_i = y_i \wedge \hat{\mathbf{c}}_i \neq \mathbf{c}_i\}
\end{equation}

La métrica propuesta, denominada \textit{Reasoning Shortcut Ratio} (RSR), se define como:
\begin{equation}
RSR = \frac{|\mathcal{A}_{RS}|}{|\mathcal{A}|}
\end{equation}

Esta métrica mide la proporción de predicciones correctas que no están sustentadas en conceptos correctos, proporcionando una estimación directa del uso de shortcuts por parte del modelo.

Adicionalmente, proponemos una variante ponderada que considera la ambigüedad estructural del espacio de conceptos. Inspirados en el análisis de óptimos deterministas, introducimos un factor de peso basado en el tamaño del conjunto $\mathcal{C}_{y_i}$:
\begin{equation}
RSR_w = \frac{\sum_{i \in \mathcal{A}_{RS}} \log \frac{1}{|\mathcal{C}_{y_i}|}}{\sum_{i \in \mathcal{A}} \log \frac{1}{|\mathcal{C}_{y_i}|}}
\end{equation}

En esta formulación, el término $\log |\mathcal{C}_{y_i}|^{-1}$ asigna mayor peso a instancias con baja ambigüedad estructural (i.e., cuando existen pocas configuraciones de conceptos que explican la etiqueta), y menor peso a aquellas con alta ambigüedad. De esta manera, la métrica prioriza la detección de shortcuts en escenarios donde estos son más críticos desde el punto de vista semántico.

Finalmente, la métrica propuesta permite: $i)$ detectar el uso de reasoning shortcuts en modelos NeSy. $ii)$ Comparar modelos más allá de métricas tradicionales como accuracy. $iii)$ Analizar el impacto de la ambigüedad conceptual en el aprendizaje.


\subsubsection{Extensión probabilística}

Hasta ahora, las métricas propuestas consideran únicamente decisiones puntuales del modelo. Sin embargo, muchos modelos NeSy modernos producen distribuciones de probabilidad sobre los conceptos. Para capturar esta información, extendemos la métrica a un marco probabilístico.

Sea un modelo que define una distribución $p_\theta(c \mid x)$ sobre el espacio de conceptos dado una entrada $x$. Nos enfocamos en instancias donde el modelo predice correctamente la etiqueta, es decir, $\hat{y}_i = y_i$. En este caso, evaluamos en qué medida el modelo asigna probabilidad al concepto correcto $c_i$. Definimos el grado de shortcut por instancia como:
\begin{equation}
RS^{\text{prob}}_i = 1 - p_\theta(c_i \mid x_i)
\end{equation}

Esta cantidad mide el \textit{desajsute probabilístico} entre el modelo y el concepto correcto. Un valor alto indica que, a pesar de predecir correctamente la etiqueta, el modelo asigna baja probabilidad al concepto verdadero, lo que sugiere el uso de explicaciones alternativas (shortcuts). Por el contrario, valores cercanos a cero indican que la predicción correcta está respaldada por una representación conceptual coherente.

A partir de esta definición local, proponemos una métrica global que agrega este comportamiento sobre todas las predicciones correctas, incorporando además la ambigüedad estructural del problema:
\begin{equation}
D^{\text{prob}}_{\text{modelo}} =
\sum_{i: \hat{y}_i = y_i}
\big(1 - p_\theta(c_i \mid x_i)\big)\,\log (|R(y_i)| -1)
\end{equation}

Esta métrica cuantifica el \textit{grado esperado de uso de shortcuts} del modelo. Cada término pondera dos factores complementarios:
\begin{enumerate}
    \item \textbf{Intensidad del shortcut:} capturada por $1 - p_\theta(c_i \mid x_i)$, que refleja cuán poco probable considera el modelo al concepto correcto.
    \item \textbf{Ambigüedad estructural:} capturada por $\log (|R(y_i)| -1)$, que mide cuántas configuraciones conceptuales pueden producir la misma etiqueta.
\end{enumerate}

De esta manera, la métrica asigna mayor peso a shortcuts en regiones del espacio donde existen múltiples explicaciones posibles, alineándose con el análisis teórico de óptimos deterministas.

En conjunto, esta extensión permite ir más allá de la detección binaria de shortcuts, capturando su \textit{intensidad} y proporcionando una caracterización más fina y continua del comportamiento del modelo.

\subsection{Escenarios Experimentales}

En el examen ENEM, los ensayos son evaluados a partir de cinco \textit{criterios}. En este trabajo nos enfocamos en el primero de ellos, C1, el cual evalúa el dominio de la modalidad escrita formal del idioma. Este criterio se descompone en nueve subcriterios organizados en dos dimensiones: (i) \textit{syntax}, con cinco niveles (non-existing, deficitary, regular, good, excellent), y (ii) \textit{mistakes}, con cuatro niveles (many mistakes, some mistakes, few mistakes, maximum 2 mistakes). Estas dimensiones son codificadas como variables ordinales, donde \textit{syntax} toma valores en $\{0, \dots, 4\}$ y \textit{mistakes} en $\{0, \dots, 3\}$.

El puntaje final de C1 se obtiene mediante un conjunto de reglas lógicas derivadas directamente de la rúbrica oficial del ENEM, las cuales definen una función determinista $K$ que asigna un \textit{score} en función de ambas dimensiones:

\[
K = \left\{
\begin{array}{ll}
\text{syntax}(0) & \rightarrow \text{score}(0) \\
\text{syntax}(1) \land \text{mistakes}(0) & \rightarrow \text{score}(1) \\
\text{syntax}(1) \land \text{mistakes}(\geq 1) & \rightarrow \text{score}(2) \\
\text{syntax}(\geq 2) \land \text{mistakes}(0) & \rightarrow \text{score}(2) \\
\text{syntax}(2) \land \text{mistakes}(\geq 1) & \rightarrow \text{score}(3) \\
\text{syntax}(\geq 3) \land \text{mistakes}(1) & \rightarrow \text{score}(3) \\
\text{syntax}(3) \land \text{mistakes}(\geq 2) & \rightarrow \text{score}(4) \\
\text{syntax}(\geq 4) \land \text{mistakes}(2) & \rightarrow \text{score}(4) \\
\text{syntax}(4) \land \text{mistakes}(3) & \rightarrow \text{score}(5)
\end{array}
\right.
\]

\subsubsection{Datasets}
%JBCS \cite{propor2024} e Subcriteria \cite{feedbackpropor}
Utilizamos el \textit{Subcriteria dataset} \cite{feedbackpropor}, el cual fue construido específicamente para modelar los subcriterios definidos en la rúbrica del ENEM. Este dataset se compone de dos fuentes: \textit{Source A}, que contiene 63 ensayos anotados manualmente por el experto A con sus respectivos subcriterios, y \textit{Source B}, que incluye los ensayos evaluados por el experto B bajo el mismo esquema. Ambas particiones contienen tanto las anotaciones a nivel de subcriterio como las calificaciones finales.

Adicionalmente, empleamos el \textit{AES ENEM dataset} \cite{propor2024}, que contiene 770 ensayos con sus calificaciones a nivel de \textit{criterio} y nota final, pero sin anotaciones explícitas de subcriterios.

\subsubsection{Modelos Evaluados}
%Distant learning on Subcriteria, supervised learning on Subcriteria
Consideramos el sistema de \textit{Automatic Essay Scoring} propuesto en \cite{feedbackpropor}, el cual sigue un enfoque neuro-simbólico basado en \textit{Scallop}. En este sistema, modelos de lenguaje (LLMs) son utilizados para predecir los subcriterios a partir de los ensayos, mientras que las reglas lógicas de la rúbrica se utilizan para inferir el puntaje final.

Esta arquitectura permite definir distintos escenarios experimentales:

\begin{itemize}
    \item \textbf{Distant learning:} utilizamos dos subredes preentrenadas (una para cada dimensión: \textit{syntax} y \textit{mistakes}) siguiendo \cite{feedbackpropor}. Estas predicciones son luego combinadas mediante el programa lógico en Scallop para obtener el score de C1. La evaluación se realiza sobre el Subcriteria dataset.

    \item \textbf{Semi-supervised learning:} dado que las anotaciones de subcriterios son limitadas, entrenamos dos modelos \textit{BERTimbau} para predecir las dimensiones latentes utilizando el AES ENEM dataset. En ausencia de anotaciones explícitas de subcriterios, incorporamos las restricciones de las reglas lógicas como una penalización en la función de pérdida. Este escenario permite analizar la estrategia de mitigación \textit{concept supervision} propuesta en \cite{Marconato2023-NotAll}.

    \item \textbf{Pretrained learning:} utilizamos dos subredes preentrenadas propuestas en \cite{feedbackpropor}, estas redes fueron entrenadas usando LLMs para predecir de forma independiente las dimensiones \textit{syntax} y \textit{mistakes}. La evaluación se realiza sobre el Subcriteria dataset.
\end{itemize}

\subsection{Métricas de evaluación}

Evaluamos los modelos considerando múltiples niveles de análisis. En primer lugar, reportamos la \textit{accuracy} de la predicción final del score de C1. En segundo lugar, utilizamos matrices de confusión a nivel de conceptos, que permiten analizar la calidad de las predicciones intermedias y detectar patrones asociados a \textit{reasoning shortcuts} \cite{Bortolotti2024}.

Finalmente, incorporamos nuestra métrica propuesta, \textit{RSR ponderada}, junto con su extensión probabilística, lo que nos permite cuantificar el grado esperado en el que un modelo recurre a shortcuts durante el proceso de inferencia.

\section{Resultados y Discusión} \label{sec:results}

En esta sección presentamos los resultados obtenidos de nuestros experimentos.


\begin{center}
\begin{tabular}{lccc}
\toprule
Escenário & Accuracy & $RSR_w$ & $RSR_{prob}$ \\
\midrule
Distant learning & 0.6 & 0.30 & 0.52 \\
Semi-supervised learning & 0.7 & 0.28 & 0.50 \\
\bottomrule
\end{tabular}
\end{center}


\section{Conclusiones} \label{sec:conclusion}

\paragraph{Resumen de Contribuciones}
\paragraph{Limitaciones}
\paragraph{Trabajo Futuro}

\section{Referencias}

\bibliographystyle{sbc}
\bibliography{sbc-template}

\end{document}
